\section{Preliminaries}
\label{sec:preliminaries}

This section introduces the foundational concepts and notation used throughout the paper.
Table~\ref{tab:notation} summarizes the key symbols.

\begin{table}[t]
\centering
\caption{Summary of notation.}
\label{tab:notation}
\begin{tabular}{cl}
\toprule
\textbf{Symbol} & \textbf{Description} \\
\midrule
$S_i$ & Stakeholder $i$ \\
$a_i^{v}$ & Artifact $i$, version $v$ \\
$\mathcal{S}_{\text{auth}}(a)$ & Set of stakeholders authorized for artifact $a$ \\
$\mathcal{T}_{\text{approved}}$ & Set of approved toolchains \\
$\pi$ & Zero-knowledge proof \\
$H(\cdot)$ & Collision-resistant hash function \\
$c_i^{v} = H(a_i^{v})$ & On-chain commitment to artifact $a_i^{v}$ \\
$\textsf{Prove}(\mathsf{pk}, x, w)$ & Proof generation algorithm \\
$\textsf{Verify}(\mathsf{vk}, x, \pi)$ & Proof verification algorithm \\
$P(\cdot)$ & Property predicate \\
$T$ & Approved tool ($T \in \mathcal{T}_{\text{approved}}$) \\
$\tau$ & Transformation identifier \\
$f_\tau$ & Deterministic transformation function \\
$\Delta$ & Version delta \\
$\boldsymbol{\pi}$ & Proof bundle (concatenation of predicate proofs) \\
\bottomrule
\end{tabular}
\end{table}

\subsection{Zero-Knowledge Succinct Non-interactive Arguments of Knowledge}
\label{sec:prelim-zksnark}

A zero-knowledge succinct non-interactive argument of knowledge (zk-SNARK) is a proof system that allows a prover to convince a verifier that a statement $x$ is true with respect to a witness $w$, without revealing $w$ and without interaction beyond a single message~\cite{groth2016size}.
Formally, a zk-SNARK for an NP relation $\mathcal{R}$ consists of three algorithms:

\begin{itemize}
    \item $(\mathsf{pk}, \mathsf{vk}) \leftarrow \textsf{Setup}(1^\lambda, \mathcal{R})$: A one-time setup that generates a proving key $\mathsf{pk}$ and verification key $\mathsf{vk}$ for relation $\mathcal{R}$, parameterized by a security parameter $\lambda$.
    \item $\pi \leftarrow \textsf{Prove}(\mathsf{pk}, x, w)$: The prover generates a proof $\pi$ that $(x, w) \in \mathcal{R}$.
    \item $\{0, 1\} \leftarrow \textsf{Verify}(\mathsf{vk}, x, \pi)$: The verifier accepts (1) or rejects (0) the proof.
\end{itemize}

A zk-SNARK satisfies three properties: \emph{completeness} (an honest prover always convinces the verifier), \emph{soundness} (no polynomial-time adversary can produce a valid proof for a false statement, except with negligible probability), and \emph{zero-knowledge} (the proof reveals nothing about the witness beyond the truth of the statement).

We use the Groth16 construction~\cite{groth2016size}, which produces constant-size proofs (2 $\mathbb{G}_1$ + 1 $\mathbb{G}_2$ elements, approximately 128 bytes compressed on BN254, or 256 bytes uncompressed) and enables verification via a single multi-pairing check over 4 input pairs.
BN254 provides approximately 100--110 bits of security following advances in the number field sieve for pairing-friendly curves; this is sufficient for the proof-of-concept, and production deployments requiring 128-bit security should consider BLS12-381 (with correspondingly larger proof sizes of approximately 192 bytes compressed).
Groth16 requires a \emph{trusted setup} ceremony that generates a structured reference string (SRS); if the toxic waste from this ceremony is not properly destroyed, soundness can be violated.
We discuss alternatives to trusted setup (e.g., PLONK, STARKs) in Section~\ref{sec:discussion}.

We denote by $H$ a collision-resistant hash function used to compute artifact commitments: $c_i^{v} = H(a_i^{v})$.
All on-chain references to artifacts use these commitments rather than the raw content, preserving confidentiality.
In our framework, the witness $w$ is the engineering artifact content, the statement $x$ is a property predicate (e.g., ``this artifact conforms to schema $S$''), and the proof $\pi$ is included in the assurance passport credential.

\subsection{Consortium Blockchain and Consensus Protocols}
\label{sec:prelim-blockchain}

A consortium blockchain is a permissioned distributed ledger where a predefined set of validator nodes maintains consensus on the shared state.
Unlike public blockchains, consortium chains restrict participation to authorized entities, providing higher throughput and finality guarantees at the cost of decentralization.

We use GoQuorum~\cite{wells2023hotstuff_goquorum}, an enterprise-grade Ethereum fork originally developed by J.P. Morgan and later maintained by ConsenSys, which supports private transactions and multiple consensus protocols.
While ConsenSys has since shifted focus to Hyperledger Besu as the primary enterprise Ethereum client, GoQuorum remains functional and was selected for our proof-of-concept due to its mature Raft support, EVM compatibility, and existing deployment on our edge hardware testbed.
Our evaluation compares three consensus mechanisms:

\textbf{Istanbul Byzantine Fault Tolerance (IBFT).}
A PBFT-inspired protocol where a designated proposer broadcasts a block and validators exchange three rounds of messages (pre-prepare, prepare, commit) to reach agreement.
IBFT tolerates up to $f = \lfloor (n-1)/3 \rfloor$ Byzantine faults among $n$ validators and provides immediate transaction finality.

\textbf{QBFT.}
A refined variant of IBFT with improved liveness guarantees and a formally verified specification, while maintaining the same fault tolerance bound.
QBFT is the recommended BFT protocol for GoQuorum production deployments.

\textbf{Raft.}
A crash fault tolerant (CFT) protocol based on leader election.
Raft assumes nodes may crash but do not behave maliciously, tolerating up to $f = \lfloor (n-1)/2 \rfloor$ crash faults.
It offers faster block times than BFT protocols but provides weaker security guarantees, making it suitable for trusted single-organization deployments.

For DME environments with mutually distrusting stakeholders, IBFT or QBFT is required; Raft serves as a performance upper bound in our benchmarks.

\subsection{W3C Verifiable Credentials}
\label{sec:prelim-vc}

The W3C Verifiable Credentials Data Model~\cite{w3c_vc} defines a standard for expressing tamper-evident credentials on the web.
A verifiable credential (VC) consists of three components:

\begin{itemize}
    \item \textbf{Credential metadata}: issuer identity, issuance date, expiration, credential type.
    \item \textbf{Claims}: a set of subject-property-value assertions (e.g., ``artifact $a$ conforms to schema $S$'').
    \item \textbf{Proof}: a cryptographic proof binding the claims to the issuer (e.g., a digital signature or a zero-knowledge proof).
\end{itemize}

A \emph{VC profile} is a domain-specific specialization of the VC data model that defines the credential schema, claim types, and proof mechanisms for a particular use case.
Examples include Open Badges 3.0 for educational achievements~\cite{open_badges3} and the UNTP Digital Product Passport for supply chain sustainability~\cite{untp_dpp}.
The assurance passport defined in Section~\ref{sec:framework} is a VC profile for engineering artifact attestation.

Credentials are issued by an \emph{issuer} (in our case, the stakeholder's ZK proof system), held by a \emph{holder} (the stakeholder), and verified by a \emph{verifier} (any ASOT participant or auditor) without contacting the issuer.
When combined with zero-knowledge proofs, the credential's proof field enables \emph{predicate attestation}: the holder can prove that specific properties hold without revealing the full credential or underlying data.
(This is distinct from traditional VC selective disclosure, which reveals a subset of claim values; predicate attestation reveals nothing beyond the truth of the predicate.)

\subsection{Digital Mission Engineering}
\label{sec:prelim-dme}

Digital Mission Engineering (DME) is the application of digital engineering principles to mission-level analysis and decision-making~\cite{dod_de_strategy, dod_me_guide2023}.
In a DME ecosystem, stakeholders use Model-Based Systems Engineering (MBSE) tools to create, maintain, and share engineering artifacts---including system architecture models (e.g., SysML, Unified Architecture Framework (UAF)), simulation configurations, requirements documents, and trade study results.

The \emph{Authoritative Source of Truth} (ASOT) is a governed repository that serves as the single reference point for the current state of all engineering artifacts.
In multi-stakeholder environments, maintaining ASOT integrity requires:
\begin{itemize}
    \item \textbf{Provenance}: tracing each artifact to its originating stakeholder and tool.
    \item \textbf{Version control}: tracking artifact evolution with immutable history.
    \item \textbf{Access governance}: controlling who can contribute, modify, or view specific artifacts.
    \item \textbf{Configuration management}: establishing and auditing baselines of approved artifact sets at decision points.
    \item \textbf{Interoperability}: supporting diverse toolchains across organizational boundaries.
\end{itemize}

In practice, many DME programs operate with a \emph{federated} ASOT: each stakeholder maintains their own authoritative source, and the system-level ASOT is a governed aggregation.
The on-chain commitment architecture described in this paper naturally supports both centralized and federated ASOT models, as each stakeholder commits independently and the blockchain serves as the federation layer.
Current ASOT implementations address these requirements through centralized governance and access control policies, but do not provide cryptographic guarantees---the gap our framework addresses.
